\documentclass[11pt,a4paper]{article}
\usepackage[utf8]{inputenc}
\usepackage[T1]{fontenc}
\usepackage[ngerman]{babel}
\usepackage{amsmath,amssymb}
\usepackage{graphicx,booktabs,xcolor}
\usepackage[margin=25mm]{geometry}
\usepackage{tcolorbox,hyperref}

\definecolor{darkblue}{HTML}{16213e}
\definecolor{accentred}{HTML}{e94560}
\definecolor{keygreen}{HTML}{2e7d32}

\title{\textcolor{darkblue}{Dok. 0016\\Obertonmoden der Basszunge: Profilierung und Inharmonizität}}
\author{Harmonikabau --- Technische Dokumentation}
\date{}

\begin{document}
\maketitle

\begin{tcolorbox}[colframe=red!60!black, colback=red!5!white]
\textbf{Hinweis:} Dieses Dokument bietet ein Erklärungsmodell.
Die FEM-Berechnung erfasst die Biegemoden eines idealisierten Balkens ---
Strömungskopplung, nichtlineare Amplitude und Spalteffekte sind nicht enthalten.
\end{tcolorbox}

\noindent\textit{\small Berechnungsskripte: Die vollständigen FEM-Berechnungen mit allen Tabellen
sind in zwei separaten Python-Skripten dokumentiert, die im selben Verzeichnis
auf GitHub abgelegt sind:
\textbf{pruefskript\_0016\_obertonmoden.py} (Profile, Gewichte, Steifigkeit, Knotenpositionen)
und \textbf{pruefskript\_0016\_chromatisch.py} (chromatische Analyse über 2~Oktaven,
$f_2/f_1$ = const bei gleichmäßigem Profil, $f_2$ absolut vs.\ $f_H$-Bereiche)
und \textbf{pruefskript\_0016\_praxis.py} (Praxis-Profilierung mit wandernder
dünnster Stelle, Verdünnungstiefe, Optimum-Suche).}

\noindent\textit{FEM-Berechnung der Biegemoden einer Basszunge (L = 70\,mm) mit
sechs verschiedenen Dickenprofilen. Modenformen, Frequenzverhältnisse,
Inharmonizität und Konsequenzen für Klang und Kammerkopplung.
Referenz: Dok.~0005--0009, 0012, 0015.}

\section{Obertöne eines Cantilevers sind nicht harmonisch}

Die Stimmzunge ist ein einseitig eingespannter Balken (Cantilever).
Die Eigenfrequenzen der Biegemoden folgen nicht der harmonischen Reihe
$(1:2:3:4\ldots)$, sondern der Cantilever-Reihe:

\begin{equation}
f_n / f_1 = (\beta_n / \beta_1)^2
\end{equation}

mit $\beta_1 L = 1{,}8751$, $\beta_2 L = 4{,}6941$, $\beta_3 L = 7{,}8548$,
$\beta_4 L = 10{,}9955$. Das ergibt die Verhältnisse $1 : 6{,}27 : 17{,}55 : 34{,}39$ ---
der 2.~Oberton liegt nicht bei der Oktave ($2\times$), sondern bei mehr als
zweieinhalb Oktaven ($6{,}27\times$).

\begin{tcolorbox}[colframe=keygreen, colback=green!5!white]
\textbf{Die Harmonischen im Klang stammen nicht von diesen Biegemoden.}
Sie entstehen durch die periodische Unterbrechung des Luftstroms --- die Zunge
wirkt als Impulsgenerator (Dok.~0013, 0015). Die Impulsform erzeugt harmonische
Obertöne $(1f, 2f, 3f\ldots)$, auch wenn die mechanischen Moden ganz woanders liegen.
\end{tcolorbox}

Warum sind die mechanischen Moden trotzdem wichtig? Weil sie bestimmen,
bei welchen Frequenzen die Zunge \textbf{resonant angeregt} werden kann ---
durch die Kammer (Dok.~0009), durch Nachbarzungen oder durch transiente
Druckschwankungen beim Spielen.

\section{Sechs Dickenprofile}

Die Profilierung verändert die Masseverteilung und die lokale Biegesteifigkeit.
Sechs Profile werden verglichen, alle mit $t_0 = 0{,}40$\,mm am Fuß.

\begin{center}
\begin{tabular}{lccccc}
\toprule
Profil & $t_\mathrm{Spitze}$ & Verdünnung & $f_1$ & $f_2/f_1$ & $f_3/f_1$ \\
\midrule
Gleichmäßig & 0,40\,mm & 0\,\% & 66,8\,Hz & 6,27 & 17,55 \\
Linear 80\,\% & 0,32\,mm & 20\,\% & 68,5\,Hz & 5,71 & 15,57 \\
Linear 50\,\% & 0,20\,mm & 50\,\% & 72,6\,Hz & 4,79 & 12,36 \\
Linear 30\,\% & 0,12\,mm & 70\,\% & 77,5\,Hz & 4,07 & 9,94 \\
Parabolisch 50\,\% & 0,27\,mm & 50\,\% & 77,9\,Hz & 5,01 & 12,89 \\
Umgek. parabolisch & 0,20\,mm & 50\,\% & 61,8\,Hz & 4,97 & 13,00 \\
\bottomrule
\end{tabular}
\end{center}

\section{Obertonverhältnisse im Vergleich}

Kein Profil kommt in die Nähe der harmonischen Reihe. Selbst bei 70\,\%
Verdünnung liegt Mode~2 bei $4{,}07\times$ statt $2\times$. Die Inharmonizität
ist eine fundamentale Eigenschaft des Cantilevers --- nur ein Saitenoszillator
hätte harmonische Obertöne.

\section{Stimmgewichte: Der gegenteilige Effekt}

Die Profilierung senkt $f_2/f_1$ --- die Moden rücken zusammen.
Stimmgewichte an der Spitze wirken \textbf{gegenteilig}: Sie senken $f_1$ stark,
aber die höheren Moden weniger. $f_2/f_1$ \textbf{steigt} mit zunehmendem Gewicht.

\begin{center}
\begin{tabular}{lccccc}
\toprule
Gewicht & Masse & \% $m_Z$ & $f_1$ & $f_2/f_1$ & $f_3/f_1$ \\
\midrule
Ohne & 0\,g & 0\,\% & 66,8\,Hz & 6,27 & 17,55 \\
0,10\,g & 0,10\,g & 5,7\,\% & 60,2\,Hz & 6,37 & 18,04 \\
0,20\,g & 0,20\,g & 11,4\,\% & 55,2\,Hz & 6,58 & 18,96 \\
0,40\,g & 0,40\,g & 22,9\,\% & 48,1\,Hz & 7,11 & 21,03 \\
1,50\,g & 1,50\,g & 85,9\,\% & 31,4\,Hz & 9,89 & 30,84 \\
\bottomrule
\end{tabular}
\end{center}

\begin{tcolorbox}[colframe=keygreen, colback=green!5!white]
\textbf{Der Grund:} Das Gewicht sitzt am Schwingungsmaximum von Mode~1.
Mode~1 wird maximal gebremst. Mode~2 hat an der Spitze weniger Amplitude
$\to$ wird weniger gebremst $\to$ das Verhältnis steigt.
\end{tcolorbox}

\section{Kombination: Profilierung + Gewicht}

In der Praxis wird beides kombiniert. Verdünnung senkt $f_2/f_1$,
Gewicht hebt es wieder an. Linear~30\,\% allein: $f_2/f_1 = 4{,}07$
bei $f_1 = 77{,}5$~Hz. Linear~30\,\% + 0,20~g: $f_2/f_1 = 4{,}64$
bei $f_1 = 50{,}2$~Hz.

\section{Gleiches $f_1$-Ziel, verschiedene Wege}

Bei $f_1 = 50$~Hz: Gleichmäßig + Gewicht ergibt $f_2/f_1 \approx 6{,}7$.
Linear~50\,\% + Gewicht: $\approx 5{,}3$. Linear~30\,\% + Gewicht:
$\approx 4{,}6$. Verdünnung gewinnt --- auch mit Gewicht.

\begin{tcolorbox}[colframe=keygreen, colback=green!5!white]
\textbf{Ansprache:} Die profilierte Zunge hat bei gleichem $f_1$
niedrigere Steifigkeit $\to$ bessere Ansprache (Dok.~0012).
Profilierung + wenig Gewicht ist besser als viel Gewicht allein.
\end{tcolorbox}

\section{Modenformen und Knotenpositionen}

Die Profilierung verschiebt die Schwingungsknoten zur Spitze hin.
Bei der gleichmäßigen Zunge liegt der erste Knoten von Mode~2 bei $x \approx 55\,\%$.
Bei Linear~30\,\% wandert er auf $x \approx 65\,\%$. Bei stark profilierten
Zungen konzentriert sich die Schwingungsenergie am leichten Ende.

\section{Inharmonizität in Cent}

Alle Obertöne liegen tausende Cent über der harmonischen Reihe.
Die Profilierung reduziert die Abweichung um 25--35\,\%.

\begin{tcolorbox}[colframe=keygreen, colback=green!5!white]
\textbf{Konsequenz:} Die Biegemoden fallen nie auf harmonische Obertöne
des Grundtons. Die Kammer-Kopplung (Dok.~0009) betrifft die mechanischen
Moden --- ob diese in der Nähe von $f_H$ liegen.
\end{tcolorbox}

\section{Steifigkeit und Grundfrequenz}

Profilierung hat zwei gleichzeitige Effekte: $f_1$ steigt (weniger Masse),
$k$ sinkt (dünnerer Querschnitt). Eine profilierte Zunge für denselben Ton
muss länger oder am Fuß dicker sein --- und hat bei gleicher Frequenz
niedrigere Steifigkeit, also bessere Ansprache (Dok.~0012).

\textbf{Das ist der eigentliche Praxisnutzen:} Niedrigere Steifigkeit bei
gleicher Grundfrequenz $\to$ bessere Ansprache.

\section{Kammer-Kopplung: Wo fallen die Moden hin?}

Mode~2 liegt bei der gleichmäßigen Zunge bei $\approx 420$\,Hz --- mitten
im typischen $f_H$-Bereich. Profilierung verschiebt Mode~2 nach unten
und kann ihn gezielt aus dem $f_H$-Bereich herausschieben.

\section{Praxis-Profilierung: Position der dünnsten Stelle}

In der Praxis ist die Zunge am dicksten bei der Niete und verjüngt sich glatt.
Die dünnste Stelle kann an der Spitze oder zwischen Spitze und Mitte liegen ---
der Verlauf ist nie abrupt.

\textbf{Überraschendes Ergebnis:} Die Position der dünnsten Stelle hat wenig
Einfluss auf $f_2/f_1$. Zwischen $x_\mathrm{min} = 0{,}7L$ und $1{,}0L$ ist
das Plateau praktisch flach. Die \textbf{Tiefe} der Verdünnung bestimmt
$f_2/f_1$, die Position bestimmt $f_1$.

\section{Optimale Verdünnungstiefe}

$f_2/f_1$ hat ein \textbf{Minimum bei $\approx 70\,\%$ Verdünnung}
($t_\mathrm{min} \approx 0{,}12$\,mm bei $t_0 = 0{,}40$\,mm). Bei stärkerer
Verdünnung steigt $f_2/f_1$ wieder an.

\begin{center}
\begin{tabular}{lcccc}
\toprule
Verdünnung & $t_\mathrm{min}$ & $f_1$ & $f_2/f_1$ & $f_3/f_1$ \\
\midrule
0\,\% & 0,40\,mm & 66,8\,Hz & 6,27 & 17,55 \\
40\,\% & 0,24\,mm & 68,4\,Hz & 4,70 & 12,61 \\
50\,\% & 0,20\,mm & 67,7\,Hz & 4,35 & 11,50 \\
60\,\% & 0,16\,mm & 65,4\,Hz & 4,07 & 10,50 \\
70\,\% & 0,12\,mm & 60,4\,Hz & 3,94 & 9,75 \\
80\,\% & 0,08\,mm & 49,8\,Hz & 4,25 & 9,61 \\
\bottomrule
\end{tabular}
\end{center}

\section{Chromatische Analyse: 2 Oktaven, gleiche Mensur}

Bei einer gleichmäßigen Zunge ist $f_2/f_1 = 6{,}27$ = \textbf{konstant},
unabhängig von der Dicke. Da $f \propto t$, skalieren alle Moden gleich ---
das Verhältnis fällt raus. Die Mensur legt $f_2/f_1$ für alle Töne fest.

Die untere Oktave (A1--A2) ist die \textbf{kritische Zone}: $f_2$ liegt bei
345--689\,Hz --- mitten im typischen $f_H$-Bereich. Die obere Oktave: $f_2 > 700$\,Hz,
weit über $f_H$. Das erklärt, warum Geistertöne hauptsächlich bei tiefen Tönen auftreten.

\section{Zusammenfassung}

\begin{enumerate}
\item \textbf{Cantilever-Obertöne sind inharmonisch:}
  $1 : 6{,}27 : 17{,}55 : 34{,}39$. Die hörbaren Harmonischen stammen
  vom Impulsgenerator, nicht von den Biegemoden.
\item \textbf{Profilierung komprimiert die Verhältnisse:} Optimum bei
  $\approx 70\,\%$ Verdünnung ($f_2/f_1 \approx 3{,}94$). Darüber hinaus
  steigt $f_2/f_1$ wieder.
\item \textbf{Position der dünnsten Stelle:} Zwischen 70\,\% und 100\,\% der
  Länge kaum Unterschied. Die Tiefe bestimmt $f_2/f_1$, die Position
  bestimmt $f_1$.
\item \textbf{Stimmgewichte spreizen:} Gegenteiliger Effekt zur Verdünnung.
  Profilierung + wenig Gewicht ist besser als viel Gewicht allein.
\item \textbf{$f_2/f_1$ = const bei gleichmäßigem Profil:} Die Mensur
  legt das Verhältnis für alle chromatischen Töne fest.
\item \textbf{Kritische Zone = untere Oktave:} $f_2$ liegt für A1--A2
  bei 345--689\,Hz --- mitten im $f_H$-Bereich. Dort treten Geistertöne auf.
\item \textbf{Hauptnutzen:} Niedrigere Steifigkeit bei gleicher $f_1$
  $\to$ bessere Ansprache. Verschiebung von Mode~2 aus dem $f_H$-Bereich.
\end{enumerate}

\bigskip
\textit{Die Zunge schwingt inharmonisch. Der Klang ist harmonisch.
Die Profilierung ändert die Ansprache --- die Obertöne macht die Luft.}

\end{document}
